% Options for packages loaded elsewhere
% Options for packages loaded elsewhere
\PassOptionsToPackage{unicode}{hyperref}
\PassOptionsToPackage{hyphens}{url}
\PassOptionsToPackage{dvipsnames,svgnames,x11names}{xcolor}
%
\documentclass[
  letterpaper,
  DIV=11,
  numbers=noendperiod]{scrartcl}
\usepackage{xcolor}
\usepackage{amsmath,amssymb}
\setcounter{secnumdepth}{-\maxdimen} % remove section numbering
\usepackage{iftex}
\ifPDFTeX
  \usepackage[T1]{fontenc}
  \usepackage[utf8]{inputenc}
  \usepackage{textcomp} % provide euro and other symbols
\else % if luatex or xetex
  \usepackage{unicode-math} % this also loads fontspec
  \defaultfontfeatures{Scale=MatchLowercase}
  \defaultfontfeatures[\rmfamily]{Ligatures=TeX,Scale=1}
\fi
\usepackage{lmodern}
\ifPDFTeX\else
  % xetex/luatex font selection
\fi
% Use upquote if available, for straight quotes in verbatim environments
\IfFileExists{upquote.sty}{\usepackage{upquote}}{}
\IfFileExists{microtype.sty}{% use microtype if available
  \usepackage[]{microtype}
  \UseMicrotypeSet[protrusion]{basicmath} % disable protrusion for tt fonts
}{}
\makeatletter
\@ifundefined{KOMAClassName}{% if non-KOMA class
  \IfFileExists{parskip.sty}{%
    \usepackage{parskip}
  }{% else
    \setlength{\parindent}{0pt}
    \setlength{\parskip}{6pt plus 2pt minus 1pt}}
}{% if KOMA class
  \KOMAoptions{parskip=half}}
\makeatother
% Make \paragraph and \subparagraph free-standing
\makeatletter
\ifx\paragraph\undefined\else
  \let\oldparagraph\paragraph
  \renewcommand{\paragraph}{
    \@ifstar
      \xxxParagraphStar
      \xxxParagraphNoStar
  }
  \newcommand{\xxxParagraphStar}[1]{\oldparagraph*{#1}\mbox{}}
  \newcommand{\xxxParagraphNoStar}[1]{\oldparagraph{#1}\mbox{}}
\fi
\ifx\subparagraph\undefined\else
  \let\oldsubparagraph\subparagraph
  \renewcommand{\subparagraph}{
    \@ifstar
      \xxxSubParagraphStar
      \xxxSubParagraphNoStar
  }
  \newcommand{\xxxSubParagraphStar}[1]{\oldsubparagraph*{#1}\mbox{}}
  \newcommand{\xxxSubParagraphNoStar}[1]{\oldsubparagraph{#1}\mbox{}}
\fi
\makeatother


\usepackage{longtable,booktabs,array}
\usepackage{calc} % for calculating minipage widths
% Correct order of tables after \paragraph or \subparagraph
\usepackage{etoolbox}
\makeatletter
\patchcmd\longtable{\par}{\if@noskipsec\mbox{}\fi\par}{}{}
\makeatother
% Allow footnotes in longtable head/foot
\IfFileExists{footnotehyper.sty}{\usepackage{footnotehyper}}{\usepackage{footnote}}
\makesavenoteenv{longtable}
\usepackage{graphicx}
\makeatletter
\newsavebox\pandoc@box
\newcommand*\pandocbounded[1]{% scales image to fit in text height/width
  \sbox\pandoc@box{#1}%
  \Gscale@div\@tempa{\textheight}{\dimexpr\ht\pandoc@box+\dp\pandoc@box\relax}%
  \Gscale@div\@tempb{\linewidth}{\wd\pandoc@box}%
  \ifdim\@tempb\p@<\@tempa\p@\let\@tempa\@tempb\fi% select the smaller of both
  \ifdim\@tempa\p@<\p@\scalebox{\@tempa}{\usebox\pandoc@box}%
  \else\usebox{\pandoc@box}%
  \fi%
}
% Set default figure placement to htbp
\def\fps@figure{htbp}
\makeatother


% definitions for citeproc citations
\NewDocumentCommand\citeproctext{}{}
\NewDocumentCommand\citeproc{mm}{%
  \begingroup\def\citeproctext{#2}\cite{#1}\endgroup}
\makeatletter
 % allow citations to break across lines
 \let\@cite@ofmt\@firstofone
 % avoid brackets around text for \cite:
 \def\@biblabel#1{}
 \def\@cite#1#2{{#1\if@tempswa , #2\fi}}
\makeatother
\newlength{\cslhangindent}
\setlength{\cslhangindent}{1.5em}
\newlength{\csllabelwidth}
\setlength{\csllabelwidth}{3em}
\newenvironment{CSLReferences}[2] % #1 hanging-indent, #2 entry-spacing
 {\begin{list}{}{%
  \setlength{\itemindent}{0pt}
  \setlength{\leftmargin}{0pt}
  \setlength{\parsep}{0pt}
  % turn on hanging indent if param 1 is 1
  \ifodd #1
   \setlength{\leftmargin}{\cslhangindent}
   \setlength{\itemindent}{-1\cslhangindent}
  \fi
  % set entry spacing
  \setlength{\itemsep}{#2\baselineskip}}}
 {\end{list}}
\usepackage{calc}
\newcommand{\CSLBlock}[1]{\hfill\break\parbox[t]{\linewidth}{\strut\ignorespaces#1\strut}}
\newcommand{\CSLLeftMargin}[1]{\parbox[t]{\csllabelwidth}{\strut#1\strut}}
\newcommand{\CSLRightInline}[1]{\parbox[t]{\linewidth - \csllabelwidth}{\strut#1\strut}}
\newcommand{\CSLIndent}[1]{\hspace{\cslhangindent}#1}



\setlength{\emergencystretch}{3em} % prevent overfull lines

\providecommand{\tightlist}{%
  \setlength{\itemsep}{0pt}\setlength{\parskip}{0pt}}



 


\KOMAoption{captions}{tableheading}
\makeatletter
\@ifpackageloaded{caption}{}{\usepackage{caption}}
\AtBeginDocument{%
\ifdefined\contentsname
  \renewcommand*\contentsname{Table of contents}
\else
  \newcommand\contentsname{Table of contents}
\fi
\ifdefined\listfigurename
  \renewcommand*\listfigurename{List of Figures}
\else
  \newcommand\listfigurename{List of Figures}
\fi
\ifdefined\listtablename
  \renewcommand*\listtablename{List of Tables}
\else
  \newcommand\listtablename{List of Tables}
\fi
\ifdefined\figurename
  \renewcommand*\figurename{Figure}
\else
  \newcommand\figurename{Figure}
\fi
\ifdefined\tablename
  \renewcommand*\tablename{Table}
\else
  \newcommand\tablename{Table}
\fi
}
\@ifpackageloaded{float}{}{\usepackage{float}}
\floatstyle{ruled}
\@ifundefined{c@chapter}{\newfloat{codelisting}{h}{lop}}{\newfloat{codelisting}{h}{lop}[chapter]}
\floatname{codelisting}{Listing}
\newcommand*\listoflistings{\listof{codelisting}{List of Listings}}
\makeatother
\makeatletter
\makeatother
\makeatletter
\@ifpackageloaded{caption}{}{\usepackage{caption}}
\@ifpackageloaded{subcaption}{}{\usepackage{subcaption}}
\makeatother
\usepackage{bookmark}
\IfFileExists{xurl.sty}{\usepackage{xurl}}{} % add URL line breaks if available
\urlstyle{same}
\hypersetup{
  pdftitle={Report: Grand Challenges in IDT - Coordinating Research and Development on Significant Problems of Practice},
  pdfkeywords={instructional design and technology, grand challenges,
problem-centered research, human-centered education, collaborative
inquiry},
  colorlinks=true,
  linkcolor={blue},
  filecolor={Maroon},
  citecolor={Blue},
  urlcolor={Blue},
  pdfcreator={LaTeX via pandoc}}


\title{Report: Grand Challenges in IDT - Coordinating Research and
Development on Significant Problems of Practice}
\author{Stephanie Moore \and Theresa Huff \and Natalie
Milman \and Matthew Schmidt \and Jason McDonald}
\date{}
\begin{document}
\maketitle
\begin{abstract}
This white paper reports on the OTESSA Santa Fe Colloquium, a two-day
scholarly convening held in June 2025 as part of the OTESSA 2025
Congress. Designed as a pilot for a collaborative, problem-centered
model of inquiry, the colloquium brought together researchers,
instructional designers, faculty, and journal editors to shift attention
from educational technologies toward significant problems of practice in
instructional design and technology (IDT). Through structured
activities, including a World Café dialogue model and a Grand Challenges
workshop, participants collectively identified, refined, and prioritized
pressing challenges facing the field. Three priority grand challenges
emerged: considering the whole human in education, supporting learner
mental wellness, and redefining the purposes and structures of
education. This paper synthesizes key insights from conference
activities and post-colloquium reflection papers, highlighting shared
themes, areas of alignment, and opportunities for coordination. Rather
than presenting empirical findings, the report offers a
practice-oriented synthesis intended to inform strategic planning,
future convenings, and coordinated research and practice efforts within
OTESSA and the broader IDT community.
\end{abstract}


\section{\texorpdfstring{\textbf{Introduction}}{Introduction}}\label{introduction}

This special issue is the result of the OTESSA Santa Fe Colloquium that
took place June 5--6, 2025, as part of the broader OTESSA 2025 Congress,
which included both virtual and in-person gatherings across multiple
sites. While the main Congress ran virtually from June 2--6 and included
on-site meetings in Toronto and Victoria, the Santa Fe event served as a
dedicated, in-person scholarly colloquium hosted at the Museum of Indian
Arts \& Culture in Santa Fe, New Mexico.

The Santa Fe convening was organized by Stephanie Moore (University of
New Mexico), Matthew Schmidt (University of Georgia), and Jason McDonald
(Brigham Young University). It brought together sixteen invited
participants, selected through a competitive call for proposals.
Attendees represented a mix of faculty, instructional designers,
researchers, and doctoral students from institutions including the
University of New Mexico, University of Georgia, Brigham Young
University, University of Minnesota, The George Washington University,
Loyola Marymount University, University of Idaho, and Texas Tech
University. Several journal editors also participated, representing
\emph{Contemporary Issues in Technology and Teacher Education},
\emph{OTESSA Journal}, \emph{Journal of Computing in Higher Education},
and \emph{Journal of Applied Instructional Design}.

This was the first convening of its kind for OTESSA, designed to pilot a
model of scholarly colloquia centered on Grand Challenges in
instructional design and technology (IDT). Inspired by similar
initiatives in other fields, the colloquium's purpose was to shift the
focus of IDT research and practice away from specific technologies and
toward complex, global problems that require coordinated,
interdisciplinary approaches.

\section{Format and Activities}\label{format-and-activities}

The two-day program combined plenary sessions, workshops, and
collaborative ``World Café'' discussions. The World Café model is a
structured small group dialogue process in which participants rotate
among tables for timed rounds of discussion, each centered on a focused
question or artifact. In each round, a table host briefly orients new
participants, after which the group engages in collaborative meaning
making, building on insights accumulated across prior rounds. This
approach leverages distributed expertise, promotes cumulative sense
making, and creates conditions for generative critique that is both
dialogic and action oriented (Hokanson 2023).

Participants were expected to arrive well-prepared, having read and
reviewed all accepted draft papers in advance. Each participant brought
questions, feedback, and ideas for how papers might connect to broader
grand challenges in the field.

\begin{itemize}
\item
  Day One opened with a keynote from Brad Hokanson (University of
  Minnesota) on the importance of finding and framing problems in design
  research. The day also featured the first Grand Challenges Hands-on
  Workshop, led by Drs. Moore, Milman, and Huff, which guided
  participants through structured brainstorming and prioritization
  exercises. Using sticky notes, colored stickers, and group huddles,
  attendees identified and refined possible grand challenges in IDT. In
  the afternoon, participants engaged in World Café Sessions,
  small-group roundtables where authors presented their work, received
  structured peer feedback, and collectively identified themes, which
  were later synthesized on a shared Miro board.
\item
  Day Two began with a live-streamed Presidential Address from OTESSA
  leaders Valerie Irvine and Stephanie Moore, situating the colloquium's
  work within the broader international mission of OTESSA. The day also
  included a Publication Pathways session with OTESSA Journal editors,
  additional World Café rounds, and a continuation of the Grand
  Challenges Workshop, where groups elaborated on top-rated problems by
  answering guiding questions such as: \emph{What is the problem? Who is
  affected? What has been tried? What evidence is needed? What networks
  or policies might help?}
\end{itemize}

Both days featured catered lunches at the Museum Hill Café, which
extended opportunities for informal dialogue and networking. The
colloquium concluded with an action-planning roundtable and synthesis
session, where participants identified opportunities for collaboration
and committed to next steps, including preparing a 500-word reflection
paper and, for presenting authors, revising their manuscripts for
possible publication in the \emph{OTESSA Journal}.

\section{Significance of the
Convening}\label{significance-of-the-convening}

By design, the Santa Fe colloquium moved beyond traditional conference
presentations to emphasize small-group, interactive, problem-centered
dialogue. The structured World Café and workshop models enabled
participants to both critique and contribute to one another's work while
collaboratively framing the grand challenges that will guide future
OTESSA initiatives. Participants also provided written reflections and
feedback, ensuring that the outcomes of this first colloquium will
inform the design of future convenings.

The Santa Fe event demonstrated the value of blending rigorous scholarly
exchange with collaborative, strategic planning. It set the stage for
what OTESSA envisions as an ongoing initiative: building networks of
researchers and practitioners who can collectively address the grand
challenges facing instructional design and educational technology in the
years ahead.

\section{Identified Grand Challenges}\label{identified-grand-challenges}

Through a structured workshop process, participants identified an
initial list of 15 potential grand challenges in instructional design
and technology. These included: demonstrating education; learning and
healthcare; forgiving; human--technology partnerships;
anti-technorationalism; human(e) education; inclusive education;
redefining education; linguistic justice; mental health; rapidly
adaptive technology learning; equal/equitable opportunity and closing
the digital divide; quality of instruction in K--12; considering the
whole student rather than only aspects that serve institutional agendas;
and strategies, guidelines, and pedagogy for the responsible use of AI.

Following group review and voting, participants consolidated these into
three priority grand challenges for deeper discussion:

\begin{enumerate}
\def\labelenumi{\arabic{enumi}.}
\item
  \textbf{Considering the Whole Human}

  \begin{itemize}
  \item
    \emph{Problem identified}: ``Currently education requires you to fit
    to education instead of education fitting to you. We need to go
    beyond intellectual capital and address the whole human.''
  \item
    \emph{Why it matters}: Participants highlighted that learners bring
    with them mental health needs, disabilities, international status,
    family responsibilities, cultural and linguistic differences, and
    other realities often overlooked by institutions. Without addressing
    these, education risks neglecting students as full humans.
  \end{itemize}
\item
  \textbf{Mental Wellness of Learners}

  \begin{itemize}
  \item
    \emph{Problem identified}: ``The mental health crisis affects
    everything including learning, perseverance, and suicide prevention.
    Faculty, staff, and designers are not trained to support this.''
  \item
    \emph{Why it matters}: Thriving learners require environments that
    attend to wellness as a foundation for learning. Participants noted
    that without explicit attention to this challenge, both individual
    learners and the broader educational system will continue to
    decline.
  \end{itemize}
\item
  \textbf{Redefining Education}

  \begin{itemize}
  \item
    \emph{Problem identified}: ``Education is too focused on
    solutionism, technology, and production.''
  \item
    \emph{Why it matters}: Current systems emphasize serving the economy
    rather than helping learners flourish. Addressing this challenge
    means questioning assumptions about education's purpose and shifting
    toward human-centered, flourishing-focused models.
  \end{itemize}
\end{enumerate}

These challenges were captured on large sticky notes during huddle
sessions and represent draft phrasing to be refined over time. After the
group reconvened, the challenges were reviewed, narrowed, and
consolidated through rounds of voting. Participants emphasized that
focusing on \emph{problems}---rather than technologies or
solutions---was a necessary shift for the field. While connections to
frameworks such as the UNESCO Sustainable Development Goals ({``{UNESCO}
and {Sustainable Development Goals} \textbar{} {UNESCO}''} n.d.) were
left implicit, many of the accepted papers explicitly drew such links.

The final output of the colloquium was a preliminary list of identified
grand challenges, with shared notes on potential avenues for
coordination and collaboration. This list was intended as a starting
point. Participants were then tasked with writing a 500-word reflection
paper to further clarify and refine these challenges in the weeks
following the event.

\section{Key Findings from Conference
Activities}\label{key-findings-from-conference-activities}

The Santa Fe colloquium confirmed that format matters: structure and
intentional design can transform scholarly exchange into strategic
collaboration. The World Café model was repeatedly highlighted as an
invaluable experience, creating the conditions for rigorous yet
empathetic dialogue that deepened both the critique of individual papers
and the collective framing of grand challenges. Participants emphasized
that this structure, including its small group and interactive
organization, allowed them to move past defensiveness and surface more
candid, actionable insights.

\subsection{Human-Centered Engagement}\label{human-centered-engagement}

One of the strongest findings was the power of human-centered design in
academic convenings. Participants described the experience as
intellectually energizing and emotionally grounding---rare in typical
conference environments where it seems the emphasis is on checking the
box to add a presentation to one's cv and not meaningful interactions,
collaboration, critique, or deep sharing of scholarship. This reinforced
a broader theme of the colloquium: if education must be redefined to
consider the whole human, so too must scholarly gatherings.

\subsection{Strategic Value of Grand Challenges
Framing}\label{strategic-value-of-grand-challenges-framing}

The integration of grand challenges elevated discussions from
project-specific details to system-level priorities, encouraging
participants to think in terms of long-term, cross-contextual impact.
While some sought more clarity about what qualifies as a grand
challenge, most agreed that the framework was essential for aligning
research with pressing global needs and for situating work within
international agendas such as the UNESCO Sustainable Development Goals.

\subsection{Positive Signals}\label{positive-signals}

Participants saw Santa Fe as a model for:

\begin{itemize}
\item
  Deep over broad: engaging more meaningfully with fewer papers rather
  than skimming many.
\item
  Listening as practice: structuring dialogue so authors first absorb
  before responding.
\item
  Constructive critique: cultivating critique and different perspectives
  was welcome, including between well-established and new, emerging
  scholars.
\item
  Balancing inclusiveness and care: creating space for welcoming all
  participants within a collegial environment that was not hierarchical,
  unlike most other academic meetings.
\item
  Network-building: generating momentum for collaborations across
  institutions, journals, and contexts, including early career scholars
  to be.
\end{itemize}

\subsection{Challenges to Address}\label{challenges-to-address}

At the same time, participants identified key barriers to effectiveness
that should be strategically addressed in future convenings:

\begin{itemize}
\item
  Infrastructure: poor audio quality and room layout undermined hybrid
  participation.
\item
  Reach: participants wanted mechanisms to ensure broader exposure to
  all papers.
\item
  Continuity: clearer post-event pathways for collaboration are needed
  to maintain momentum.
\end{itemize}

\section{Implications for OTESSA}\label{implications-for-otessa}

These findings suggest that OTESSA can claim a distinctive space by
redefining what academic convenings look like. By designing events that
are structured, problem-centered, and human-focused, OTESSA not only
surfaces research but also builds collective capacity to address grand
challenges. At the same time, addressing practical barriers (technical,
structural, logistical) will be crucial for scaling this model and
ensuring equitable participation across diverse sites and modalities.

\subsection{Key Findings from the Participants'
Papers}\label{key-findings-from-the-participants-papers}

\subsubsection{Importance of Grand
Challenges}\label{importance-of-grand-challenges}

Across the participants' 500-word reflection papers, authors endorsed
grand challenges as a clarifying way to coordinate efforts on issues
that are bigger than any single course or tool. Contributors repeatedly
framed the problems as systemic, implicating policy, incentives, and the
purposes of education, while also cautioning that challenges must be
scoped tightly enough to act on. As one writer put it, they left with
``a sense of raised awareness,'' alongside an ``alarm'' that demanded
sharper focus; another noted that the frame ``raised questions about the
very purpose and design of education itself.''

\subsubsection{Specific Challenges
Developed}\label{specific-challenges-developed}

Across the reflection papers, three grand challenge areas surfaced most
often: reimagining education's aims beyond compliance, considering the
whole human in learning design, and supporting mental health and
wellbeing, especially in online contexts. Participants often linked
these concerns to crisis responsive and systemic design. They described
converging threads across groups: ``different approaches that
nonetheless converged around the need to attend to the whole student
(belonging, mental health, anxiety), and systemic or crisis responsive
design,'' and a shared shortlist emerging in the workshop that included
``Wellness of the Online Learner'' and ``Redefining Education.''

\subsubsection{Connection to Own Work}\label{connection-to-own-work}

Moving from priorities to practice, many authors described immediate
shifts in how they position their projects, from reframing questions to
committing to specific next steps. Several expressed a desire to align
ongoing studies with one of the identified grand challenges and to
continue reflecting on the implications for method and audience. One
participant wrote, ``I feel an urge to better situate my work within a
grand challenge\ldots{} not only from a pragmatic perspective but also
intellectually honest,'' while another reported already moving to
application: ``At my institution I've already begun\ldots{} to situate
everyday decisions within a broader social and educational context.''

\subsubsection{Influence on Paper
Thinking}\label{influence-on-paper-thinking}

Extending those adjustments into writing, the colloquium clarified how
papers connect to grand challenges, generated new ideas, and exposed
productive gaps, with several noting helpful, actionable feedback. For
some, the effect was confirmatory; for others, it prompted substantive
reframing and scope shifts beyond the course level. As one contributor
explained, the conversations ``pushed me to\ldots{} revise the paper so
that it also considers the whole faculty,'' and another noted that their
work ``recognize{[}s{]} learners as whole human beings, not just as
people engaging with content or technology.''

\subsubsection{Colloquium Feedback}\label{colloquium-feedback}

Finally, participants credited the event's design with enabling rapid
sense making and momentum. The World Café surfaced shared patterns
efficiently; plenaries and presentations provided conceptual anchors;
preparation helped people arrive ready; and Miro harvesting made
collective thinking visible. Setting also mattered, with the Santa Fe
setting described as conducive to deep dialogue. One author wrote, ``The
world café model peer discussion helped surface common themes\ldots,''
while another highlighted that there was ``something powerful about
being able to watch in real time as reflections and thematic patterns
emerged visually on the board.'' Several offered practical ideas for
future colloquia to extend post event continuity and time for deeper
dives.

\section{Overlapping and Reinforced
Findings}\label{overlapping-and-reinforced-findings}

Taken together, the live activities and the written submissions tell a
coherent story. What surfaced in the room largely reappears on the page,
with participants naming similar priorities, using similar language, and
moving from ideas to concrete next steps. Some overlapping findings
include:

\begin{enumerate}
\def\labelenumi{\arabic{enumi}.}
\item
  \textbf{Systemic lens, shared purpose.}\\
  In both settings, people looked beyond individual classes to the
  larger problems shaping education. They reported ``a sense of raised
  awareness'' and ``alarm,'' and pressed on what education is for and
  how it is designed. This shared framing points to system-level
  responses rather than isolated interventions.
\item
  \textbf{Whole human and well-being as a throughline.} Building on
  that, both in-person conversations and reflection papers centered
  belonging, anxiety, and care as essential to meaningful learning. As
  one summary put it, groups ``converged around the need to attend to
  the whole student (belonging, mental health, anxiety), and systemic or
  crisis responsive design.'' The papers then carried this emphasis
  forward, especially for online learners.
\item
  \textbf{Event structure leading to reframing and application.}\\
  The format seemed to do more than collect ideas. It prompted shifts
  that the writing then documented. Authors reported being internally
  ``pushed\ldots{} to revise the paper'' and expressed ``an urge to
  better situate my work within a grand challenge,'' with some already
  taking steps to ``situate everyday decisions within a broader social
  and educational context.''
\item
  \textbf{Rapid synthesis reinforced by visible harvesting.} Rotating
  dialogue paired with shared capture tools produced quick pattern
  recognition that persisted into the manuscripts. Participants noted
  that ``the world café model peer discussion helped surface common
  themes,'' and that it was ``powerful'' to watch patterns emerge ``in
  real time.'' The written submissions reflect those same clusters.
\item
  \textbf{Anchoring concepts enabling depth.}\\
  Plenary ideas and presentations provided conceptual scaffolds that
  showed up in the papers as tools that clarified alignment, new
  connections, and a widened scope. One participant appreciated a
  posture of being ``curious and open rather than trying to control
  every outcome,'' a stance mirrored in the manuscripts.
\item
  \textbf{Preparation and setting amplifying contributions.}\\
  Finally, participants credited pre-colloquium reading work and the
  Santa Fe setting with enabling focus and candid exchange. That
  preparation is visible in the manuscripts as clearer alignment
  statements, more specific problem definitions, and tighter rationales
  for why a grand-challenge lens fits their work.
\end{enumerate}

In sum, the live activities catalyzed synthesis and momentum, and the
papers captured that momentum as shared language, clarified purpose, and
actionable pathways for continued collaboration.

\section{Strategic Implications \& Next
Steps}\label{strategic-implications-next-steps}

This section turns the shared findings into concrete moves the community
can take now.

\begin{enumerate}
\def\labelenumi{\arabic{enumi}.}
\item
  \textbf{Publish toward the grand challenges}\\
  Journal editors and conference chairs should issue themed calls that
  require a short alignment statement, so new work advances the selected
  challenges.
\item
  \textbf{Secure a planning grant}\\
  Apply for a small planning grant that funds coordination, cross-site
  convenings, a shared repository, pilot seeding, and light evaluation.
\item
  \textbf{Build a cross-sector, international community of practice}\\
  Create a network that links K--12, community colleges, and
  universities with adjacent fields to inform instruction, policy,
  teaching, and infrastructure.
\item
  \textbf{Run small pilots that create usable evidence}\\
  Launch short, scoped pilots tied to specific challenges and share
  brief case write-ups and artifacts others can adopt.
\item
  \textbf{Provide simple tools people can use tomorrow}\\
  Release a one-paragraph alignment box for manuscripts and proposals
  and a short well-being and belonging checklist for courses and
  programs.
\item
  \textbf{Connect research to leadership decisions}\\
  Prepare a two-page brief for deans and program heads that states why
  the challenges matter now and asks for one concrete support action.
\item
  \textbf{Track a few shared indicators}\\
  Agree on three easy measures such as pilots launched and completed,
  courses using the checklist, and a short belonging or anxiety item.
\item
  \textbf{Keep an open channel for critique}\\
  Invite counterpoints and refinements through a simple form and publish
  responses so the challenge set remains focused and useful.
\item
  \textbf{Map scholarship to international frameworks}\\
  Conduct a systematic literature review that maps existing working to
  international frameworks such as UNESCO's SDGs. Encourage authors and
  journals to ``tag'' an SDG by incorporating relevant SDGs into the
  article's keywords list.
\end{enumerate}

\section{Conclusion}\label{conclusion}

This first OTESSA colloquium accomplished what was hoped. It brought
people together to name a focused set of grand challenges, link those
challenges to ongoing research and work, and use thoughtful event
structures that can move ideas into action. Across the live sessions and
the papers, participants focused on problems rather than solutions, the
need to support the whole human, and clear ways to carry the work
forward.

The path ahead is to keep attention on the challenges, seek planning
funds to support coordination where possible, and continue building a
cross-sector community that includes K--12, community colleges,
universities, and relevant partner fields. Publishing toward the
challenges and aligning future convenings will help sustain momentum.

We invite continued dialogue through publications, future convenings,
and shared resources that speak directly to these challenges. This white
paper is a reference point for that ongoing conversation and a guide for
practical next steps.

\section{Author's Contributions}\label{authors-contributions}

S.M., M.S., and J.M. co-conceived and co-designed the OTESSA Santa Fe
Colloquium, with S.M. providing overall leadership for the convening.
S.M., N.M., and T.H. co-led the Grand Challenges workshop; M.S. led and
coordinated the World Café Sessions and coordinated the synthesis of
colloquium activities and participant reflections. All authors
contributed to facilitation of conference sessions; T.H. analyzed
reflections for themes and strategic framing of findings and drafted the
manuscript; All authors reviewed, revised, and approved the final
version. All authors have read and approved the final manuscript.

\section{Open Researcher and Contributor Identifier
(ORCID)}\label{open-researcher-and-contributor-identifier-orcid}

\section{Ethics Statement}\label{ethics-statement}

This paper reports on a scholarly convening and synthesizes themes from
collaborative discussions and voluntary post-conference reflections. No
human subjects research was conducted, and no institutional ethics
approval was required.

\section{Conflict of Interest}\label{conflict-of-interest}

The authors declare no conflicts of interest.

\section{Data Availability Statement}\label{data-availability-statement}

No datasets were generated or analyzed as part of this practice-oriented
work.

\section*{References}\label{references}
\addcontentsline{toc}{section}{References}

\phantomsection\label{refs}
\begin{CSLReferences}{1}{0}
\bibitem[\citeproctext]{ref-hokansonAECTSummerResearch2023}
Hokanson, Brad. 2023. {``The {AECT} Summer Research Symposium.''} In
\emph{{AECT} at 100: {A} Legacy of Leadership}, edited by C. T. Miller,
A. A. Piña, M. H. Molenda, P. L. Harris, and B. B. Lockee, 589--98.
Brill.

\bibitem[\citeproctext]{ref-UNESCOSustainableDevelopment}
{``{UNESCO} and {Sustainable Development Goals} \textbar{} {UNESCO}.''}
n.d. Accessed December 17, 2025. \url{https://www.unesco.org/en/sdgs}.

\end{CSLReferences}




\end{document}
